\subsection{Генерация обобщенных z-векторов и формирование выборок}

Для эффективного захвата нелокальных во времени корреляций, присущих хаотическим системам,
 мы используем концепцию обобщенных z-векторов, построенных на основе предопределенных шаблонов.

Шаблоном называется целочисленный вектор $\boldsymbol{k} = (k_1, k_2, \dots, k_{L-1})$,
 где $L$ — длина z-вектора (число используемых наблюдений), а $1 \le k_j \le K_{\max}$ — шаг между $j$-м и $(j+1)$-м наблюдением в векторе.

Для каждого временного ряда $Y^{(m)} \in \mathcal{D}_{\text{train}}$ и каждого шаблона
 $\boldsymbol{k}$ из множества всех комбинаторно возможных шаблонов $\aleph(L, K_{\max})$
  генерируется собственная выборка обобщенных z-векторов. Z-вектор в момент времени $t$ определяется как:
\[
\mathbf{z}_{t}^{(m, \boldsymbol{k})} = \bigl(y^{(m)}_t,\ y^{(m)}_{t+k_1},\ y^{(m)}_{t+k_1+k_2},\ \dots,\ y^{(m)}_{t + \sum_{j=1}^{L-1} k_j}\bigr).
\]


Ключевым шагом нашего подхода является \textbf{объединение} всех выборок, сгенерированных
 по одному и тому же шаблону $\boldsymbol{k}$, но из разных временных рядов, в единое множество:
\[
\mathcal{Z}^{(\boldsymbol{k})} = \bigcup_{m=0}^{M} \left\{ \mathbf{z}_{t}^{(m, \boldsymbol{k})} \right\}_{t}.
\]
Таким образом, для каждого шаблона $\boldsymbol{k}$ формируется одна, но значительно
 более плотная и разнообразная обучающая выборка $\mathcal{Z}^{(\boldsymbol{k})}$,
  аккумулирующая информацию о динамике всего семейства систем.

\subsection{Кластеризация и формирование библиотеки мотивов}

Каждая объединенная выборка $\mathcal{Z}^{(\boldsymbol{k})}$ кластеризуется отдельно
 для выявления повторяющихся паттернов поведения. Центры полученных кластеров
  интерпретируются как \textit{мотивы}.
\textbf{Мотивом}, соответствующим шаблону $\boldsymbol{k}$,
 называется вектор $C_{\boldsymbol{k}} = (\eta_{\boldsymbol{k},1},
 \eta_{\boldsymbol{k},2}, \dots, \eta_{\boldsymbol{k},L})$, являющийся
 центром кластера в пространстве z-векторов $\mathcal{Z}^{(\boldsymbol{k})}$.

Множество всех мотивов, полученных для шаблона $\boldsymbol{k}$,
обозначается $\Psi_{\boldsymbol{k}}$. Полная \textit{библиотека мотивов}
 $\Psi$ представляет собой множество пар (шаблон,мотивы) для всех возможных
  шаблонов и определяется следующим образом:
\[
\Psi = \left\{ (\boldsymbol{k}, \Psi_{\boldsymbol{k}}) \mid \boldsymbol{k}
\in \aleph(L, K_{\max}) \right\}.
\]
Эта библиотека кодирует обобщенную динамику, свойственную всему классу рассматриваемых
 систем, и служит основой для построения прогнозов.

\subsection{Формирование множества возможных предсказаний}

Прогнозирование точки $y_{t+h}^{(0)}$ целевого ряда $Y^{(0)}$
 выполняется путем поиска в библиотеке $\Psi$ мотивов, чье начало
  соответствует недавней истории целевого ряда. Процесс состоит из следующих шагов:
\begin{enumerate}
    \item Для каждого шаблона $\boldsymbol{k}$ из истории целевого
     ряда $\{y_{t}^{(0)}, y_{t-1}^{(0)}, \dots\}$ формируется \textit{вектор}
      $C_{\text{query}}^{(\boldsymbol{k})}$, состоящий из $L-1$ элементов
       и структурно идентичный усеченному мотиву.

    \item Этот вектор сравнивается со всеми усеченными мотивами $C_{\boldsymbol{k},
     \text{trunc}} = (\eta_{\boldsymbol{k},1}, \dots, \eta_{\boldsymbol{k}, L-1})$
     из множества $\Psi_{\boldsymbol{k}}$.

    \item Если евклидово расстояние $\rho(C_{\text{query}}^{(\boldsymbol{k})},
     C_{\boldsymbol{k}, \text{trunc}})$ меньше заданного порога $\varepsilon$,
      то последний элемент соответствующего полного мотива, $\eta_{\boldsymbol{k},L}$,
       считается возможным предсказанием. Множество таких предсказаний для одного шаблона
       обозначается $\hat{S}_{t+h}^{(\boldsymbol{k})}$.

    \item Итоговое \textit{множество возможных предсказаний} $\hat{S}_{t+h}$ для
     точки $y_{t+h}^{(0)}$ является объединением предсказаний, полученных по всем шаблонам:
    \[
    \hat{S}_{t+h} = \bigcup_{\boldsymbol{k} \in \aleph(L, K_{\max})} \hat{S}_{t+h}^{(\boldsymbol{k})}.
    \]
\end{enumerate}

Полученное множество $\hat{S}_{t+h}$ является эмпирическим распределением и служит
 основой для последующего анализа. Если распределение значений в $\hat{S}_{t+h}$
  является плотным и унимодальным, точка считается прогнозируемой, и для нее
  вычисляется единое прогнозное значение (UPV). Если же распределение разреженное,
   мультимодальное или равномерное, точка классифицируется как непрогнозируемая.
    Детальное описание методов классификации и вычисления UPV представлено в следующем разделе.

\section{Вычисление прогноза и идентификация непрогнозируемых точек}

Центральным этапом предлагаемого подхода является анализ эмпирического распределения
 $\hat{S}_{t+h}$, полученного для каждой прогнозируемой точки $y_{t+h}^{(0)}$. Этот
  анализ решает две взаимосвязанные задачи: определение, является ли точка прогнозируемой,
   и, если да, вычисление для нее единого прогнозного значения (UPV). В данной работе обе
   задачи
   решаются с помощью единого инструмента~ кластеризации множества $\hat{S}_{t+h}$ алгоритмом DBSCAN \cite{Aggarwal2013}. 
   В ходе экспериментов использовались параметры $\varepsilon=0.01$ и $\text{min\_samples}=4$,
\subsection{Интегрированный подход на основе кластеризации}

Для каждого непустого множества $\hat{S}_{t+h}$ выполняется процедура кластеризации,
 которая разбивает его на $N_c$ кластеров $\{Q_1, Q_2, \dots, Q_{N_c}\}$. Пусть $|Q_j|$~--- число
  элементов (мощность) $j$-го кластера. Кластеры упорядочиваются по убыванию мощности:
   $|Q_1| \ge |Q_2| \ge \dots \ge |Q_{N_c}|$.

\textbf{Идентификация непрогнозируемых точек} производится на основе анализа этой
кластерной структуры. Точка $y_{t+h}^{(0)}$ считается \textit{прогнозируемой} тогда
 и только тогда, когда в множестве ее предсказаний $\hat{S}_{t+h}$ существует один
 доминирующий кластер. Формально, это условие выражается следующим образом:
\[
\zeta(\hat{S}_{t+h}) =
\begin{cases}
1, & \text{если } |Q_1| \ge \alpha \cdot |Q_2| \text{ (для } N_c > 1\text{) или } N_c = 1, \\
0, & \text{в противном случае,}
\end{cases}
\]
где $\alpha > 1$- гиперпараметр, задающий требуемую степень доминирования крупнейшего кластера.
 Во всех экспериментах, если не указано иное, использовалось значение $\alpha = 10/3$. Если самый большой кластер не превосходит по размеру
 второй по величине кластер как минимум в $\alpha$ раз, это свидетельствует о
 наличии нескольких конкурирующих, равновероятных сценариев будущего, что делает
  прогноз ненадежным. В этом случае точка классифицируется как \textit{непрогнозируемая}.
   Если множество $\hat{S}_{t+h}$ пусто, точка также считается непрогнозируемой.

\textbf{Вычисление единого прогнозного значения (UPV)} производится только для
 точек, классифицированных как прогнозируемые ($\zeta(\hat{S}_{t+h})=1$). В качестве итогового прогноза $\hat{y}_{t+h}$ принимается среднее арифметическое значение элементов, входящих в состав крупнейшего (доминирующего) кластера $Q_1$:
\[
\hat{y}_{t+h} = \frac{1}{|Q_1|} \sum_{\hat{y} \in Q_1} \hat{y}.
\]
Такой подход позволяет отфильтровать шумовые или менее согласованные группы
предсказаний и основывать прогноз на наиболее статистически значимом сценарии.

\subsection{Алгоритм кластеризации для выявления мотивов}

Для выявления типичных, повторяющихся паттернов (мотивов) в каждой
объединенной выборке z-векторов $\mathcal{Z}^{(\boldsymbol{k})}$ применяется
алгоритм кластеризации, предложенный Лапко и Ченцовым на основе идей
Уишарта \cite{LapkoChentsov2000}.

Ключевая идея метода заключается в непараметрической оценке плотности
распределения данных в многомерном пространстве. Алгоритм идентифицирует
плотные скопления точек и определяет их границы,
что делает его эффективным для поиска структур произвольной формы
в фазовом пространстве, характерных для аттракторов динамических систем.

В результате работы алгоритма для каждой выборки $\mathcal{Z}^{(\boldsymbol{k})}$
формируется набор устойчивых кластеров. Центроиды этих кластеров
и принимаются в качестве искомых \textbf{мотивов}, которые затем используются
для построения библиотеки $\Psi$.
В нашей работе для алгоритма кластеризации использовались параметры 
k=11,$\mu$=0.2. Данный выбор соответствует параметрам, примененным в работе \cite{GromovBaranov2023}.